\section*{अध्याय \ref{ch:b3:behavioural-ecology} --- व्यवहार-पारिस्थितिकी}

\begin{solution}{exo:b3:behavioural-ecology:1}
तंत्र: वह नर्तकी भोजन तथा सूरज के बीच के उस कोण को, जो उसने उड़ान पर
नापा, छत्ते पर गुरुत्व से कोण पर उतार देती है, और उस यात्रा के दृक्
प्रवाह से दूरी को उस दौड़ की अवधि के रूप में कूटबद्ध करती है।
परिवर्धन: बड़े हिस्से में जन्मजात --- जो मधुमक्खियाँ बिना कोई नृत्य देखे
पाली जाएँ वे भी ठीक नाचती हैं, और सीखे गए छोटे समायोजनों के साथ।
प्रकार्य: वह छत्ते की साथियों को किसी लाभदायक स्रोत तक भर्ती करता है और
उस बस्ती की आवक बढ़ाता है। इतिहास: भोजन के लिए निकलने की रूढ़ की गई
इरादा-गतियों से व्युत्पन्न; बौनी मधुमक्खियाँ अब भी किसी क्षैतिज पृष्ठ पर
सीधे उस स्रोत की ओर इशारा करती हुई नाचती हैं, यानी वही पुरखा रूप।
\end{solution}

\begin{solution}{exo:b3:behavioural-ecology:2}
कोई \omterm{def:b3:behavioural-ecology:tinbergen}{नियत क्रिया-प्ररूप} किसी सरल संकेत से, यानी \omterm{def:b3:behavioural-ecology:tinbergen}{संकेत उद्दीपक} से, छोड़ा
गया कोई रूढ़ क्रम है, जो एक बार शुरू होने पर पूरा होता है: उस हंस का
अंडा लुढ़काना, जो किसी अंडे-सरीखी वस्तु से छूटता है; उस गल-चूज़े का
ठोंगना, जो किसी चोंच-सरीखे आकार पर लाल धब्बे से छूटता है। तीन लाल
पट्टियों वाली वह छड़ी, जिस पर किसी असली सिर से अधिक ठोंगा गया, ने दिखाया
कि वह छोड़ने वाला तंत्र किसी लक्षण पर सुरबद्ध कोई छन्ना है --- लाल
विपर्यास, लंबाई --- न कि माता-पिता की कोई पहचान, और कि उस लक्षण को
बढ़ा-चढ़ा देना उस अनुक्रिया को बढ़ा-चढ़ा देता है।
\end{solution}

\begin{solution}{exo:b3:behavioural-ecology:3}
किसी क्षेत्रक को तब छोड़िए जब उसमें लाभ-दर गिरकर पूरे आवास की औसत दर पर,
यात्रा समेत, आ जाए। यात्रा-काल दुगुना करना उस औसत को गिरा देता है,
इसलिए वह आहार-खोजी हर क्षेत्रक में अधिक देर रुकता है, हर एक को अधिक पूरी
तरह चुकाता है, और अंत में उसकी कुल दर नीची रहती है।
\end{solution}

\begin{solution}{exo:b3:behavioural-ecology:4}
ESS कोई ऐसी रणनीति है कि एक बार लगभग सब उसे खेलने लगें तो कोई बिरला
विकल्प उससे अधिक अंक न पा सके। सारे-कबूतर ESS नहीं है: कोई अकेला बाज़ हर
मुक़ाबला जीतता है और $V$ कमाता है, कबूतरों के $V/2$ के सामने। उस मिश्रण
में हर रणनीति $V(1 - V/C)/2$ कमाती है, जो उस $V/2$ से कम है जो
सारे-कबूतर देते, क्योंकि मुक़ाबलों का कोई अंश ऐसी लड़ाइयाँ हैं जो घायल
करती हैं; और कोई व्यष्टि उस पर अकेला सुधार नहीं कर सकता, इसलिए सब उस
हानि में फँसे रहते हैं।
\end{solution}

\begin{solution}{exo:b3:behavioural-ecology:5}
$V = 6$, $C = 10$: $p^{*} = 0.6$; हर एक को प्रतिफल $6(1 - 0.6)/2 = 1.2$,
जबकि सारे-कबूतर के लिए $3$। $V = 12 > C$: बाज़ ही शुद्ध ESS है, प्रतिफल
$(12 - 10)/2 = 1$, जबकि सारे-कबूतर के लिए $6$ --- यानी वह मुक़ाबला उस
मूल्य का पाँच-छठा भाग नष्ट कर देता है।
\end{solution}

\begin{solution}{exo:b3:behavioural-ecology:6}
$\tau = T_{0}$: $t^{*} = 1.15\times 2 = \qty{2.3}{min}$, अंश $1 -
\mathrm{e}^{-1.15} = 0.68$, $R = 0.68A/4.3 = 0.16A$ प्रति मिनट। $\tau =
4T_{0}$: $t^{*} = 2.0\times 2 = \qty{4}{min}$ (अधिक ठीक-ठीक $3.9$),
अंश $0.86$, $R = 0.86A/12 = 0.072A$ प्रति मिनट।
\end{solution}

\begin{solution}{exo:b3:behavioural-ecology:7}
सौतेले भाई-बहन $1/4$; दादा-दादी--पोता-पोती $1/4$; पहले चचेरे $1/8$।
कोई अगुणित-द्विगुणित श्रमिक और उसका भाई: वह अगुणित है और उसके पास केवल
उनकी माँ के जीन हैं; और उस श्रमिक के जीनों में से आधे उस माँ से आए और
उनमें से आधे वही प्रतियाँ हैं जो उसे मिलीं, इसलिए उसके दृष्टिकोण से
$r = 1/4$ (और उसके दृष्टिकोण से $1/2$: संबंधिता तब सममित नहीं रहती जब
गुणिता भिन्न हो)। रानी से बेटा: उसके सारे जीन उसी के हैं, पर उसके अपने
जीनों का केवल आधा उसमें है: उसकी ओर से $r = 1/2$, और उसकी ओर से $1$।
\end{solution}

\begin{solution}{exo:b3:behavioural-ecology:8}
$n\,r\,b > c$, जिसमें $b = 0.01$, $c = 0.02$: यानी $n\,r > 2$। सगे
भाई-बहन, $r
= 1/2$: $n \ge 5$। भतीजियाँ, $r = 1/4$: $n \ge 9$।
असंबंधित: कभी नहीं --- और इसीलिए बिखरते नर पुकारते नहीं।
\end{solution}

\begin{solution}{exo:b3:behavioural-ecology:9}
हर \qty{25}{cm} पर संगम दुगुने, उत्तरजीविता पाँचवाँ भाग गिरती: $L = 25$:
$0.5\times 1.2 = 0.6$; $50$: $1$; $75$: $2\times 0.8 = 1.6$; $100$:
$4\times 0.6 = 2.4$; $125$: $8\times 0.4 = 3.2$; $150$: $16\times 0.2 =
3.2$। इन शर्तों पर वह गुणनफल प्राकृतिक \qty{50}{cm} से कहीं आगे तक चढ़ता
रहता है, इसलिए उस प्रतिरूप में कुछ लागतें छूट रही हैं: उड़ान में किसी
लंबी पूँछ का उत्तरजीविता-दंड रैखिक से तीखा है, मादा की अनुक्रिया संतृप्त
हो जाती है, उन पंखों का उगना ऊर्जा लेता है, और ऐंडरसन का घोंसलों का
दुगुना होना किसी एक मौसम भर नापा गया था, जीवन भर नहीं। प्राकृतिक पूँछ
वहीं है जहाँ असली लागतें काटती हैं।
\end{solution}

\begin{solution}{exo:b3:behavioural-ecology:10}
प्रतिकारी: बाज़ के सामने $(V - C)/2$, कबूतर के सामने $V/2$, और
प्रतिकारी के सामने $V/2$। सारे-प्रतिकारी समष्टि में कोई बिरला बाज़
प्रतिकारियों से मिलता है, जो पलटकर लड़ते हैं, और $(V - C)/2 < V/2$
कमाता है: वह $V < C$ होने पर हमला नहीं कर सकता। कोई बिरला कबूतर
$V/2$ कमाता है, यानी ठीक प्रतिकारियों का प्रतिफल, इसलिए कबूतर उदासीन
हैं और बहकर भीतर आ सकते हैं; और एक बार आम हो जाने पर वे बाज़ों को भीतर
आने देते हैं (कोई बाज़ किसी कबूतर के सामने $V$ कमाता है), और उस
तीन-रणनीति क्रीड़ा का कोई सरल ESS नहीं रहता --- यानी इसका पहला उदाहरण कि
कोई विकल्प जोड़ देना किसी स्थिर मिश्रण को कैसे अस्थिर कर सकता है।
\end{solution}

\begin{solution}{exo:b3:behavioural-ecology:11}
यदि सब ठीक $t_{0}$ तक डटे रहें, तो $t_{0} +
\varepsilon$ तक डटता कोई उत्परिवर्ती नगण्य अतिरिक्त लागत पर हर मुक़ाबला
जीत लेता, और ऐसे उत्परिवर्तियों की किसी समष्टि को $t_{0} + 2\varepsilon$
हरा देता, और ऐसे ही: कोई नियत समय स्थिर नहीं, इसलिए ESS को कोई
यादृच्छिक होना ही पड़ता है। किसी चरघातांकी वितरण के साथ, आगे $s$ तक डटने की प्रायिकता, अब तक
$t$ तक डटे रहने को देखते हुए, $\mathrm{e}^{-s/m}$ है, चाहे
$t$ कुछ भी हो --- यानी वह स्मृतिहीनता गुण --- इसलिए अब भी प्रदर्शन करता
कोई प्रतिद्वंद्वी इस बारे में कुछ नहीं बताता कि वह और कितनी देर चलेगा,
और कोई रणनीति बीते समय का दोहन नहीं कर सकती।
\end{solution}

\begin{solution}{exo:b3:behavioural-ecology:12}
श्रमिक किसी बहन को $3/4$ और किसी भाई को $1/4$ मूल्य देती हैं, इसलिए
श्रमिक-नियंत्रण में उस बस्ती को मादाओं में नरों से तीन गुना निवेश करना
चाहिए, यानी निवेश से $3{:}1$ का अनुपात; और वह रानी, जो दोनों से बराबर
संबंधित है, $1{:}1$ पसंद करती है। ट्राइवर्स और हेयर (1976) ने किसी एक
ही, एक बार संगम करने वाली रानी वाली चींटियों में निवेश-अनुपात
$3{:}1$ के पास पाए: श्रमिक जीतते हैं। बहु-संगम किसी श्रमिक की अपनी बहनों
से औसत संबंधिता $1/4$ से $1/2$ की ओर गिरा देता है और श्रमिकों के इष्टतम
को $1{:}1$ की ओर धकेल देता है; और बहु-संगमी रानियों वाली जातियों की
बस्तियाँ सचमुच अधिक बराबरी से निवेश करती हैं, जैसा वह तर्क पूर्वानुमान
करता है।
\end{solution}

\begin{solution}{pb:b3:behavioural-ecology:1}
\textbf{1.} $R(t) = A(1 - \mathrm{e}^{-t/T_{0}})/(\tau + t)$; और तब
छोड़िए जब $g'(t) = R(t)$: $(A/T_{0})\mathrm{e}^{-t/T_{0}} = A(1 - \mathrm{e}^{-t/T_{0}})/
(\tau + t)$।
\textbf{2.} $(\tau + t)\mathrm{e}^{t/T_{0}}/A$ से गुणा कीजिए: $(\tau + t)/T_{0}
= \mathrm{e}^{t/T_{0}} - 1$। $\tau = T_{0}$ और $x = t/T_{0}$ के साथ: $2 + x =
\mathrm{e}^{x}$; और $x = 1.15$ देता है $3.15$, जबकि $3.16$ चाहिए।
\textbf{3.} $\tau = 4T_{0}$: $5 + x = \mathrm{e}^{x}$; $x = 1.94$ देता है
$6.94$, जबकि $6.96$; $t^{*} \approx 1.94\,T_{0}$, यानी कहिए $2.0$।
\textbf{4.} सघन: $t^{*} = \qty{3.5}{min}$, $60(1 - \mathrm{e}^{-1.15}) =
\qty{41}{kJ}$, \qty{68}{\%}, $R = 41/6.5 = \qty{6.4}{kJ/min}$। विरल:
$t^{*} = \qty{5.8}{min}$, \qty{51}{kJ}, \qty{86}{\%}, $R = 51/17.8 =
\qty{2.9}{kJ/min}$।
\textbf{5.} पाँच मिनट: $g(5) = 60(1 - \mathrm{e}^{-5/3}) = \qty{49}{kJ}$।
सघन: $49/8 = \qty{6.1}{kJ/min}$ (इष्टतम का \qty{96}{\%}); विरल:
$49/17 = \qty{2.9}{kJ/min}$ (\qty{99}{\%})। कोई कच्चा नियम कुछ ही
प्रतिशत खोता है --- यानी इतना अच्छा कि वरण उसी पर मान गया।
\textbf{6.} पकड़ों के बीच का अंतराल क्षेत्रक के चुकते ही लंबा होता जाता
है, इसलिए कोई नियत छोड़ने-का-समय कोई नियत सीमांत पकड़-दर है --- यानी उस
प्रमेय की कसौटी, जो हर क्षेत्रक में वही है। वह इसलिए विफल होता है कि
सही देहली उस आवास पर निर्भर है: नियत \qty{20}{s} के साथ वह प्राणी वहाँ
बहुत जल्दी छोड़ देता है जहाँ यात्रा लंबी है और वहाँ बहुत देर से जहाँ वह
छोटी है, जब तक वह छोड़ने-का-समय स्वयं हाल के अनुभव से न सधे --- जो
आहार-खोजी करते ही हैं।
\textbf{7.} बाज़ बनाम बाज़ $(10 - 30)/2 = -10$, बाज़ बनाम कबूतर $10$,
कबूतर बनाम बाज़ $0$, कबूतर बनाम कबूतर $5$; $p^{*} = V/C = 1/3$।
\textbf{8.} $W_{H} = 10 - 20p$, $W_{D} = 5(1 - p)$। $p = 0.1$: $8$ बनाम
$4.5$, बाज़ फैलते हैं। $p = 1/3$: हर एक $3.33$, संतुलन। $p = 0.6$: $-2$
बनाम $2$, कबूतर फैलते हैं।
\textbf{9.} ESS पर $3.33$, जबकि सारे-कबूतर के लिए $5$: यानी हर मुक़ाबले
के मूल्य का तिहाई लड़ने की संभावना में चला जाता है।
\textbf{10.} $V = 40 > C$: बाज़ ही शुद्ध ESS है; हर मुक़ाबला कोई लड़ाई
है। कठिन वर्षों में अधिक लड़ाइयाँ और अधिक चोटें दिखनी चाहिए --- और
दिखती हैं।
\textbf{11.} किसी बुर्जुआ से मिलता कोई बाज़ उसे आधी बार निवासी पाता है
(कोई लड़ाई, $(V - C)/2$) और आधी बार घुसपैठिया (कोई आसान जीत, $V$):
औसत $(3V - C)/4$। बुर्जुआ अपनों के बीच $V/2$ कमाता है। बाज़ केवल तभी
हमला कर पाता है जब $(3V - C)/4 > V/2$ हो, यानी $V > C$। उस तितली का
``निवासी जीतता है'' बुर्जुआ है: कोई सस्ती रूढ़ि, जो इसलिए स्थिर है कि
लड़ना किसी धूप-धब्बे से अधिक महँगा है --- और जब दोनों अपने को निवासी
मानें तब दोनों बाज़ खेलते हैं।
\textbf{12.} सबसे अच्छा क्या करना है, यह इस पर निर्भर है कि बाकी सब
क्या करते हैं, इसलिए वह समष्टि वहाँ जमती है जहाँ रणनीतियाँ बराबर
फ़ायदा दें, न कि किसी व्यष्टि के सर्वोत्तम पर।
\textbf{13.} $n\bar r\,b > c$: यानी $n\bar r > c/b = 3$।
\textbf{14.} आठ सगे भाई-बहन: $n\bar r = 4 > 3$, ड्यूटी फ़ायदेमंद है। आठ
अजनबी: $0$, कभी नहीं। आठ सौतेले भाई-बहन: $2 < 3$, नहीं।
\textbf{15.} $c = 0.01$: देहली $n\bar r > 1$ --- अब सौतेले भाई-बहन और
कुछ सगे भाई-बहन भी काफ़ी हैं। जिस व्यष्टि की लागत सबसे छोटी है उसी के
लिए वह नियम पूरा होता है, इसलिए अच्छी तरह खाए हुए पहरा देते हैं और भूखे
खोदते हैं।
\textbf{16.} $k$ बहनें $\tfrac{3}{4}k$ की हैं, $k$ बेटियाँ
$\tfrac{1}{2}k$ की: बहनें पसंद कीजिए, $3{:}2$ से।
\textbf{17.} कोई यादृच्छिक बहन सगी बहन ($3/4$) है प्रायिकता $1/3$ से,
और सौतेली बहन ($1/4$) प्रायिकता $2/3$ से: $\bar r = 0.25 +
0.17 = 0.42 < 0.5$। वह पसंद पलट जाती है: अब बेटियाँ बहनों से आगे निकल
जाती हैं, इसलिए अकेली संबंधिता बहुसंगमी जातियों में बंध्य श्रमिकों की
व्याख्या नहीं कर सकती --- उसमें पारिस्थितिकी और बस्ती-स्तर की चौकसी
जोड़नी पड़ती है।
\textbf{18.} उस नियम को केवल इतना चाहिए कि जिनकी मदद की जाए उनमें $r$
समष्टि से औसतन ऊँचा हो; और बंधुत्व से सहसंबद्ध कोई भी संकेत काम दे देता
है --- जन्म-बिल के पास रहना, उनकी मदद करना जिनके साथ बड़े हुए। किसी
भू-गिलहरी का नियम ``घर के पास हो तब पुकारो'' हो सकता है, जो बिना किसी को
पहचाने बंधुओं को छाँट लेता है।
\textbf{19.} $W = (L/20)(1 - L^{2}/10^{4})$: $L = 20$: $0.96$; $40$:
$1.68$; $60$: $1.92$; $80$: $1.44$।
\textbf{20.} $W' = \tfrac{1}{20}(1 - 3L^{2}/10^{4}) = 0$: $L^{*} = 100/\sqrt{3}
= \qty{58}{cm}$, $W = 2.89\times\tfrac{2}{3} = 1.92$।
\textbf{21.} $L^{*} = 70/\sqrt{3} = \qty{40}{cm}$: यानी छोटी। वह इष्टतम
वहीं बैठता है जहाँ सीमांत संगम-लाभ सीमांत उत्तरजीविता-हानि के बराबर हो,
इसलिए कोई भारी परभक्षण-लागत उस पूँछ को भीतर खींच लेती है, और भिन्न
परभक्षण के तहत समष्टियों के आभूषण भिन्न होने चाहिए --- जैसे गप्पी तथा
स्वोर्डटेल के होते हैं।
\textbf{22.} दोनों से। भगोड़ा पूर्वानुमान करता है कि वह पसंद उस लक्षण से
आगे निकल जाती है, जो उस उत्तरजीविता-इष्टतम पर रुक जाता है; और बाधा
पूर्वानुमान करती है कि मादाएँ लंबी इसलिए पसंद करती हैं कि उसे केवल कोई
तंदुरुस्त नर ढोता है, यानी फिर उससे आगे जितना अधिकांश उगा सकते हैं। वह
प्रयोग इष्टतम से आगे की कोई पसंद दिखाता है और उन हिसाबों के बीच फ़ैसला
नहीं करता।
\textbf{23.} यदि वह पूँछ मुफ़्त होती, तो हर नर सबसे लंबी उगा लेता, वह
उसकी गुणवत्ता के बारे में कुछ न कहती, और उसके लिए कोई पसंद दोहन का
शिकार होकर क्षय हो जाती। कोई ऐसी लागत, जो किसी बुरी दशा वाले नर पर
अधिक भारी पड़े, उस लंबी पूँछ को केवल तंदुरुस्तों के वहन में लाती है, और
तब वह पसंद सुयोग्यता छाँटती है: वह लागत ही उस संकेत को ईमानदार बनाती है।
\textbf{24.} प्रति संगम आठ गुने प्रतिफल के साथ, किसी नर की सुयोग्यता संगमों
के साथ तीखी चढ़ती है और किसी मादा की मुश्किल से; इसलिए नर संगमों के लिए
स्पर्धा करने को वरित होते हैं --- यानी बाज़--कबूतर मुक़ाबले तथा उनके
हथियार --- और मादाएँ, जिनकी थोड़ी संतति में से हर एक मायने रखती है,
चुनने को वरित होती हैं, जिससे नर-आभूषण फ़ायदेमंद हो जाते हैं। प्रति
संगम लाभ की कोई एक असममिति वे लड़ाइयाँ और वह सजावट, दोनों पैदा कर देती
है।
\textbf{25.} सघन आवास में $t^{*} \approx \qty{3.5}{min}$; $p^{*}
= 1/3$; और प्रहरी-ड्यूटी तब फ़ायदेमंद है जब $n\bar r > 3$।
\end{solution}
